% !TEX root = ../main.tex

% 中英标题：\chapter{中文标题}[英文标题]
\chapter{绪\hspace{1em}论}[Introduction]
% 绪论也需像摘要、目录格式那样空出两个半角字符。

\section{课题背景及研究的目的和意义}

\subsection{课题背景}

近年来，随着深度学习技术的快速发展，以大语言模型（Large Language Models, LLMs）为代表的生成式人工智能取得了突破性进展。
尤其是自 OpenAI 于 2022 年发布 ChatGPT 以来，人工智能研究进入了以大规模预训练模型为核心的新阶段。
LLM 凭借其强大的语义理解能力与文本生成能力，在自然语言处理等多个领域中展现出优异的性能~\cite{llm,ouyang2022training,llama}。
与传统任务驱动的模型不同，通用大语言模型通过在海量跨领域语料上进行预训练，具备较强的跨任务迁移能力，
为通用人工智能的发展提供了重要基础。

然而，随着研究的不断深入，LLM 的局限性也逐渐显现。首先，通用大语言模型虽然具备良好的泛化能力，
但在特定领域任务中往往难以超越专用模型，导致在实际应用中仍需针对具体场景进行定制化优化。
其次，LLM 的知识主要来源于预训练阶段所使用的数据，其知识具有明显的“静态性”，难以及时反映现实世界中的动态变化。
当面对需要最新信息或精确知识的任务时，模型容易产生过时或错误的回答。
此外，在医疗等对数据安全性要求较高的领域，受限于可用训练数据，LLM 的表现往往受到进一步制约。

在知识密集型任务（knowledge-intensive tasks）中，尤其是涉及多跳推理的问答任务，LLM 的上述问题更加突出。
模型在处理超出训练数据范围或需要跨文档推理的问题时，容易产生“幻觉”现象，
即生成与事实不符的内容~\cite{knowledge-intensive, hallucinations}。
为了解决这一问题，检索增强生成（Retrieval-Augmented Generation, RAG）方法被提出，
如\figref{fig: 对比图}所示，相比于直接将用户问题输入进行LLM生成，
RAG方式通过从外部知识库动态检索相关信息，并将其作为上下文输入模型，能够提示模型利用最新、相关的知识进行推理与生成，
从而增强生成结果的事实性与可靠性~\cite{lewis2020retrieval,rag,izacard2022few}。

\begin{figure}[htbp]
	\centering
	\includegraphics[width=0.8\textwidth]{直接生成和RAG对比图.png}
	\caption{LLM直接生成和RAG方法对比图}
	\label{fig: 对比图}
\end{figure}

近年来，RAG 已成为提升 LLM 性能的重要技术路径，并在多个研究工作中得到广泛应用。
然而，现有 RAG 方法在面对大规模非结构化语料时仍面临显著挑战。
由于相关信息通常分散在多个异构文档中，检索得到的上下文往往规模庞大且缺乏结构组织，导致噪声信息大量引入
，从而影响模型的推理准确性与生成连贯性~\cite{irrelevant-context,irrelevant-noise,asai2024self}。
尽管部分研究通过将文档划分为更细粒度的片段以提升检索效率\cite{borgeaud2022improving, jiang2023active}，
但该策略容易破坏原有语义结构，导致关键上下文信息丢失，从而降低复杂推理任务的性能~\cite{han2024retrieval,zhang2025survey}。

为缓解上述问题，研究人员进一步尝试将结构化知识引入 RAG 框架中，
其中知识图谱（Knowledge Graph, KG）~\cite{knowledge-graphs} 被广泛用于建模实体及其关系，以提升多跳推理能力。
这类方法通过显式建模结构化知识，在一定程度上提高了检索的相关性与推理的可解释性。
然而，由于知识图谱通常依赖自动化关系抽取构建，其质量容易受到噪声与错误传播的影响，进而影响整体性能。

因此，如何在降低上下文噪声的同时提升语义表达能力与推理能力，成为当前检索增强生成研究中的关键问题。
围绕这一问题，本文从结构化推理与语义建模两个角度出发，探索高效且鲁棒的检索增强生成方法。

\subsection{研究目的与意义}

针对当前大语言模型在知识密集型多跳问答任务中存在的知识静态性、推理能力不足以及检索增强生成方法中上下文噪声较大等问题，
本文旨在探索更加高效且鲁棒的检索增强生成方法，以提升模型在复杂推理场景中的表现能力。

具体而言，本文从“结构化推理”与“语义建模”两个角度出发，研究如何在引入外部知识的同时，
有效组织与表达检索到的上下文信息。一方面，通过构建显式的推理结构，引导模型逐步完成多跳推理过程，以减少无关信息干扰；
另一方面，通过引入多粒度语义表示与上下文融合机制，提升检索信息的表达能力与利用效率，从而在无需大规模模型训练的前提下，增
强大语言模型的推理性能与泛化能力。

本文的研究工作具有以下几个方面的重要意义：

（1）\textbf{提升多跳问答任务的推理能力}：通过引入结构化推理路径与语义增强机制，
有效缓解传统 RAG 方法在复杂推理任务中的能力不足问题，提升模型对跨文档、多步推理问题的处理能力。

（2）\textbf{降低上下文噪声并提高信息利用效率}：针对检索结果中存在的冗余与噪声问题，
本文通过结构引导与语义建模相结合的方式，提高上下文信息的信噪比，从而提升模型生成结果的准确性与稳定性。

（3）\textbf{提高方法的可解释性与可扩展性}：结构化推理过程能够显式展示模型的推理路径，
而联合语义建模则增强了上下文的表达能力，两者结合不仅提高了模型性能，也为结果提供了更好的可解释性。
同时，该方法无需对大语言模型进行重新训练，具有良好的扩展性与应用潜力。

（4）\textbf{降低应用成本并增强实际可用性}：相比于重新训练领域专用大模型，本文方法通过改进检索与上下文构建策略，
在保持较低计算成本的同时实现性能提升，为大语言模型在实际场景中的落地应用提供了更加可行的解决方案。

综上所述，本文的研究不仅在方法层面上对检索增强生成技术进行了有益探索，
也在实际应用层面为提升大语言模型在复杂任务中的表现提供了新的思路。

\section{国内外研究现状}

\subsection{基于外部文档的检索增强生成}

基于外部文档的检索增强生成方法主要面向非结构化文本语料，通过检索器从外部语料库中选取与查询最相关的若干文档片段，
并将其直接拼接到输入上下文中，以辅助大语言模型生成答案。
该类方法通常不显式构建图结构，也不额外设计结构化推理路径，而是依赖语义相似度检索与生成模型自身的上下文理解能力完成问答，
因此具有实现简单、易于扩展和适配现有大语言模型的优点。

从发展脉络上看，REALM\cite{guu2020retrieval} 较早将“可检索的外部知识存储”引入语言模型预训练过程，
证明了通过外部文档检索增强语言模型能够有效提升开放域问答性能。
随后，Lewis\cite{lewis2020retrieval} 等人提出的 RAG 模型进一步明确了“参数化记忆 + 非参数化记忆”的统一范式，
即利用生成模型作为参数化知识载体，同时通过向量索引的外部文档库提供可更新的显式知识支持，从而显著提高知识密集型任务的效果。
与之相关，DPR\cite{karpukhin2020dense} 则为开放域问答提供了高效的稠密检索基础，其双编码器框架在检索准确率和推理效率之间取得了良好平衡，
成为后续大量 RAG 方法的默认检索组件之一。

在此基础上，研究者从不同角度扩展了基于外部文档的 RAG 的能力边界。一类工作将检索机制更深地融入语言模型架构中，
例如 RETRO\cite{borgeaud2022improving} 通过在语言模型解码过程中引入检索到的邻近文本块，在大规模语料库上实现了显式记忆增强，
表明外部检索可以在较少参数规模下带来显著收益。另一类工作则关注少样本与上下文学习场景，
例如 Atlas\cite{izacard2022few} 通过联合训练检索器与生成模型，显著提升了知识密集型任务中的少样本学习能力；
In-Context RALM\cite{ram2023context} 则进一步提出无需改动模型结构、仅通过检索文档并直接拼接到输入提示中的方式，即可在多种场景下获得稳定收益。
这些研究共同表明，基于非结构化文本直接检索并拼接上下文的范式具有很强的通用性和工程可落地性。

尽管基于外部文档的 RAG 在开放域问答和知识增强生成中表现出良好的实用性，但其局限性也较为明显。
首先，该类方法通常依赖文档块级别的语义相似度进行检索，难以显式捕捉跨文档的多跳关系，
因此在复杂推理任务中容易遗漏关键证据。其次，将检索到的多个文本片段直接拼接到上下文中，
往往会引入大量冗余信息和噪声，增加模型的阅读负担，影响最终生成结果的准确性与稳定性。
此外，为了提高召回率，基于外部文档的 RAG 往往需要扩大检索范围或缩小切块粒度，
但这又可能带来上下文割裂、语义不完整或输入长度迅速增长等问题。
正是由于这些不足，后续研究开始尝试引入知识图谱、开放式信息抽取、多粒度语义建模等机制，
以进一步提升检索增强生成在复杂场景下的推理能力与信息利用效率。

\subsection{基于公开知识图谱的检索增强生成}

为提升检索增强生成方法在复杂推理任务中的表现，研究者开始引入结构化知识以弥补纯文本检索的不足。
其中，一类重要方法是直接利用已有的公开知识图谱作为外部知识源，
将其与大语言模型相结合，以增强模型的推理能力与知识表达能力。

公开知识图谱通常由大规模结构化三元组构成，能够显式刻画实体及其关系，为多跳推理提供天然的结构支持。
例如，在医疗领域中，SPOKE\cite{spoke} 和 UMLS\cite{umls} 等大型知识图谱被广泛应用于医学问答任务，
通过引入实体关系网络来辅助模型理解复杂的医学知识体系。
基于此类知识图谱的 RAG 方法通常将查询映射到图谱中的相关实体，并通过图结构进行扩展检索，从而获取与问题相关的多跳知识路径。

在具体方法上，已有研究从多个角度探索了知识图谱与大语言模型的融合方式。
一类方法通过构建基于知识图谱的提示信息，将图谱中的结构化知识转化为自然语言输入，从而增强模型的上下文表达能力。
例如，Soman\cite{soman2023biomedical} 等人提出的方法通过从医学知识图谱中提取相关子图，
并将其转化为提示信息，有效提升了大语言模型在医学问答任务中的表现。
类似地，KG-Rank\cite{yang2024kg} 方法进一步结合知识图谱与排序机制，通过对候选答案进行图结构约束下的重排序，
提高了答案选择的准确性。
另一类方法则尝试将图结构直接融入检索或表示学习过程中。例如，GNN-RAG\cite{mavromatis2025gnn} 利用图神经网络
（Graph Neural Network, GNN）\cite{he2021improving,mavromatis2022rearev}对知识图谱进行编码，通过建模实体之间的结构关系，实现更加精细的语义检索。
该方法能够在图结构中进行信息传播，从而更好地捕捉多跳推理所需的隐含关系，提升大语言模型的推理能力。

此外，一些通用知识图谱（如 Wikidata\cite{vrandevcic2014wikidata}、DBpedia\cite{auer2007dbpedia} 等）也被用于开放域问答任务，通过实体链接与图结构扩展，
实现从结构化知识到文本生成的映射。这类方法通常依赖实体识别与链接技术，将自然语言问题中的关键实体映射到知识图谱节点，
并通过关系路径搜索获取相关知识，从而为模型提供更加精准的上下文信息。
尽管基于公开知识图谱的方法在提升推理能力和可解释性方面具有显著优势，但其也存在一定局限性。
首先，公开知识图谱的覆盖范围有限，难以包含所有领域或最新知识，导致在开放域场景中存在知识缺失问题。
其次，知识图谱的构建与维护成本较高，且不同图谱之间存在结构差异，限制了方法的通用性。
此外，将结构化图谱知识转换为适合大语言模型输入的文本形式，仍然存在信息损失与表达不充分的问题。

综上所述，基于公开知识图谱的检索增强生成方法通过引入结构化知识，有效提升了模型的推理能力与解释性，
但其对图谱质量与覆盖范围的依赖也限制了其在复杂开放场景中的应用。
这一局限性进一步推动了后续研究从“静态知识图谱”向“动态构建知识结构”的方向发展，
即通过开放式信息抽取从文本中自动构建知识图谱，从而提升方法的灵活性与适应性。

\subsection{基于开放式信息抽取的检索增强生成}

相较于直接利用 Wikidata、UMLS、SPOKE 等公开知识图谱作为外部知识源的方法，
基于公开知识图谱的检索增强生成往往面临两个现实问题：
一方面，公开知识图谱的覆盖范围有限，难以充分适配开放域场景；
另一方面，现有公开知识图谱大多集中于医疗、生物等特定领域，通用性不足。
基于此，近年来更常见的思路是利用开放式信息抽取（Open Information Extraction, OpenIE）或大语言模型，
从原始文本中动态抽取实体关系三元组，进而构建面向当前语料库的知识图谱或知识索引，
以提升检索增强生成在复杂问答任务中的推理能力与适应性。

这类方法的核心思想在于：首先从文档集合中抽取实体、关系及其对应的语义单元，构建图结构或三元组集合；
随后，围绕用户查询在该知识结构上执行实体扩展、路径搜索、图传播或图表示学习，以获取更具推理价值的上下文信息；
最后，再将筛选后的结构化知识、对应文本证据或两者的融合结果提供给大语言模型完成答案生成。
与基于外部文档的 RAG 直接基于文本块相似度进行检索不同，基于开放式信息抽取的方法能够显式建模文本内部的关系结构，
因此在多跳问答、跨文档证据整合和复杂推理任务中显示出更强的潜力。

现有基于开放式信息抽取的检索增强生成研究大致可以分为以下几类。

第一类方法强调基于自动抽取三元组构建推理链，以提升多跳推理能力。
TRACE\cite{fang2024trace} 通过从检索文档中自动生成知识图谱，并进一步构建由多个逻辑相连三元组组成的推理链，将其作为回答问题的核心证据。
该方法说明，与直接将完整检索文档输入大语言模型相比，三元组推理链能够更有效地组织分散证据，降低噪声干扰。
第二类方法侧重于在特定专业领域中动态构建图谱，并通过图谱与文本的协同增强提升问答性能。
例如，DALK\cite{li2024dalk} 面向阿尔茨海默病科学文献场景，利用大语言模型从文献中持续抽取知识并构建演化中的领域知识图谱，
同时结合粗到细采样与自感知检索机制，实现知识图谱与大语言模型的动态协同增强。
类似地，KAG\cite{liang2025kag} 则面向专业领域知识服务，强调知识图谱与原始文本块的双向索引、逻辑形式引导的混合推理以及知识对齐机制，
以提升领域问答中的逻辑推理和专业知识利用能力。
第三类方法将开放式信息抽取得到的图结构作为图检索或图表示学习的基础，以增强检索阶段的结构感知能力。
GraphRAG\cite{graphrag} 是这一方向的代表性工作，其首先从语料中抽取实体和关系构建知识图谱，再通过社区划分与社区摘要实现从局部到全局的问答与总结。
LightRAG\cite{lightrag} 在此基础上进一步强调轻量化与效率，通过双层检索机制同时支持低层实体检索与高层关系聚合。
KET-RAG\cite{ketrag} 则从索引成本角度出发，提出多粒度索引框架，仅对关键文本块构建知识图谱骨架，
并用文本—关键词二部图作为轻量替代，以在降低构建成本的同时保持较好的检索质量。
第四类方法尝试直接利用图神经网络或图基础模型对自动构建的知识结构进行建模，以提升对复杂图关系的表达能力。
例如，GFM-RAG\cite{gfm-rag} 提出了面向 RAG 的图基础模型，通过在大规模知识图与文档上进行预训练，增强模型在未见数据集上的泛化能力；
G-Retriever\cite{g-retriever} 则关注文本图理解与问答任务，通过检索增强生成方式对图结构与文本信息进行联合建模，从而更好地支持基于图的复杂问答。
第五类方法更强调记忆建模与长程知识组织。HippoRAG\cite{hipporag} 将开放式信息抽取得到的知识图谱与个性化 PageRank 相结合，
模拟“长期记忆”式的信息检索机制，以支持多跳推理；HippoRAG2\cite{hipporag2} 在此基础上进一步引入段落节点和包含关系边，
并联合语义相似度与结构相似度，提升检索的稳定性与事实性。
与之相关，LinearRAG\cite{linearrag} 重新审视了图增强检索的流程，不再强调复杂关系图，而是通过轻量级实体抽取与语义链接构建线性化的层级图结构，
并结合局部语义桥接与全局重要性聚合，实现高效、抗噪的多跳检索。

\section{本文的主要研究内容}

总体来看，基于开放式信息抽取的检索增强生成方法兼具灵活性与结构建模能力。
相比直接依赖公开知识图谱的方法，该类方法能够从原始语料中动态构建知识结构，具有更强的领域适应性与语料可扩展性；
相比基于外部文档的朴素 RAG 方法，又能够通过显式建模实体关系，更好地组织多跳推理所需的证据，从而提升复杂问答任务中的推理能力。
然而，这一类方法仍然面临若干关键挑战：
其一，自动抽取的实体关系三元组不可避免地存在噪声与遗漏，容易引入局部语义错误，从而影响后续推理过程的可靠性；
其二，图谱构建、索引维护以及路径搜索通常带来额外计算开销，限制了方法在大规模语料场景中的效率与可扩展性；
其三，结构化知识与原始文本语义之间仍存在表达鸿沟，如何实现两者的高效融合，以提升信息利用效率与生成质量，仍是当前研究的核心问题。

针对上述问题，本文围绕“结构化推理”与“语义建模”两个关键方向，对基于开放式信息抽取的检索增强生成方法进行了系统研究，主要工作如下：

（1）针对多跳问答任务中上下文信噪比较低的问题，本文提出了一种基于三元组骨架引导的检索增强生成方法。
该方法首先利用大语言模型从检索文档中抽取实体关系三元组，构建文档级知识图谱；
随后，通过迭代扩展与评分机制生成推理骨架，将其作为结构化推理路径；
在此基础上，基于推理骨架对原始文本进行精确索引，获取与推理路径一致的文本证据，
从而构建“结构引导 + 文本补充”的上下文表示方式。该方法能够有效过滤冗余信息，显著提升上下文信噪比，
并增强模型在多跳推理任务中的表现能力。

（2）针对图结构方法依赖关系抽取且构建成本较高的问题，本文进一步提出了一种基于联合语义建模的检索增强生成方法。
该方法在 LinearRAG 框架基础上，对语义表示与上下文构建过程进行优化：
在语义抽取阶段，引入实体与短语的多粒度语义抽取机制，以增强语义覆盖能力；
在表示学习阶段，提出基于句子加权的段落表示方法，通过加权聚合句子嵌入刻画段落内部语义分布；
在上下文构建阶段，提出联合语义上下文范式，将实体、短语与段落信息进行统一组织，从而提升上下文的信息密度与表达能力。
该方法在无需显式推理路径的情况下，实现了高效且稳定的性能提升。

围绕上述方法，本文在多个多跳问答基准数据集上开展了系统实验，从性能表现、消融分析、上下文构建方式以及样例分析等多个角度，
对所提出方法的有效性进行了全面验证，并对不同方法在推理能力、信息利用效率与计算开销之间的权衡关系进行了深入分析。

综上所述，本文通过从结构化推理与语义建模两个角度对检索增强生成方法进行改进，提出了两种互补的解决方案，为提升大语言模型在复杂多跳问答任务中的性能与效率提供了新的思路。

\section{本文组织结构}

本文围绕检索增强生成方法在多跳问答任务中的语义建模与推理能力提升问题展开研究，
主要从结构化推理与多粒度语义建模两个方面进行探索。全文共分为五个章节，具体组织结构如下：

第一章为绪论。首先从大语言模型的发展背景及其在知识密集型任务中的局限性出发，阐述检索增强生成技术的研究意义；
随后对国内外相关研究现状进行分析，重点介绍朴素 RAG、基于知识图谱的 RAG 以及基于开放式信息抽取的 RAG 方法的发展情况，
并在此基础上引出本文的研究问题与主要研究内容。

第二章为相关技术概述。围绕检索增强生成系统的关键组成部分，对相关基础技术进行系统介绍，
包括检索增强生成方法、基于知识图谱的结构化建模方法、文本表示与向量检索技术，以及命名实体识别与开放式信息抽取技术，
为后续方法设计提供理论基础。

第三章深入研究基于三元组骨架引导的检索增强生成方法。
首先分析传统 RAG 在多跳推理任务中存在的上下文噪声与结构缺失问题，
随后提出通过构建知识三元组推理骨架来引导文本检索与上下文构建的方法，并详细介绍推理骨架的构建与评分机制。
最后通过多组实验验证该方法在提升推理能力与降低上下文冗余方面的有效性。

第四章进一步研究基于联合语义建模的检索增强生成方法。
在分析现有图结构方法构建成本高与语义表达不足问题的基础上，提出多粒度语义抽取、句子加权段落表示以及联合语义上下文构建方法，
实现对实体、短语与文本信息的统一建模。通过实验分析验证该方法在无需显式推理路径的情况下，能够有效提升检索与生成性能。

结论部分总结了本文在针对知识库问答的检索增强生成方法上的研究成果，分析方法的优势与局限，
并对未来在检索增强生成、多粒度语义建模及复杂推理任务中的研究方向进行展望。